\documentclass[12pt,a4paper,oneside,openany]{report}

\usepackage[utf8]{inputenc}
\usepackage[T2A]{fontenc}
\usepackage[russian]{babel}
\usepackage{cmap}
\usepackage{geometry}
\usepackage{booktabs}
\usepackage{longtable}
\usepackage{array}
\usepackage{float}
\usepackage{placeins}
\usepackage{amsmath}
\usepackage{xcolor}
\usepackage{verbatim}
\usepackage{graphicx}
\graphicspath{{images/}}
\usepackage{hyperref}

\geometry{left=30mm,right=15mm,top=20mm,bottom=20mm}
\hypersetup{
  unicode=true,
  colorlinks=true,
  linkcolor=black,
  citecolor=black,
  urlcolor=blue
}
\sloppy
\emergencystretch=2em

\usepackage{etoolbox}
\makeatletter

\patchcmd{\chapter}{\thispagestyle{plain}}{\thispagestyle{plain}\setcounter{footnote}{0}}{}{}
\makeatother


\widowpenalty=1000
\clubpenalty=1000


\hyphenpenalty=1000
\exhyphenpenalty=100



\begin{document}

\begin{titlepage}
  \centering
  \includegraphics[width=0.45\textwidth]{msu_logo.png}\\[0.5em]

  {\normalsize Московский государственный университет имени М.\,В.\,Ломоносова}\\[0.2em]
  {\normalsize Факультет космических исследований}\\[0.2em]
  {\normalsize Специальность 01.05.01 ``Фундаментальные математика и механика''}\\[0.2em]
  \noindent\rule{0.85\textwidth}{0.4pt}

  \vfill

  {\LARGE\textbf{Исследование методов применения больших нейросетевых языковых моделей\\
  для пополнения онтологии в специальной предметной области\\
  ракетно-космической техники}}\\[0.5em]
  {\normalsize Курсовая работа}

  \vfill

  \begin{flushright}
  \begin{minipage}{0.55\textwidth}
    \raggedleft
    Выполнил студент 4 курса группы 402\\
    Лисс Михаил Алексеевич\\[1em]
    Научный руководитель\\
    Добров Борис Викторович
  \end{minipage}
  \end{flushright}

  \vfill

  Москва, 2026
\end{titlepage}

\tableofcontents
\clearpage

\chapter*{Введение}
\addcontentsline{toc}{chapter}{Введение}

Онтологии и тезаурусы описывают понятия предметной области и связи между ними. В информационно-поисковом тезаурусе основная единица обычно не словоформа. Важнее понятие: к нему привязаны варианты записи термина, синонимы и отношения с другими понятиями. Для автоматической обработки текстов такая схема удобна. Например, разные написания одного термина можно привести к одной записи, а связи между понятиями затем использовать в поиске, рубрикации и пополнении ресурса \cite{louk2010,ontotez2006}.

В таксономической части онтологии одно из ключевых отношений связано с гиперонимией. Для заданного понятия нужно найти одно или несколько более общих родовых понятий. Для термина ``спутник связи'' естественными кандидатами будут ``искусственный спутник Земли'', ``космический аппарат'' или более специальный класс, если он уже есть в тезаурусе. При пополнении онтологии это принципиально: новый термин мало добавить в список, его нужно поставить в уже существующую структуру знаний \cite{gruber1993,hearst1992}.

В последние годы для пополнения таксономий начали использовать большие языковые модели. В работе Ф.\,А.~Садковского, Н.\,В.~Лукашевич и И.\,Ю.~Гришина ``Автоматическое пополнение таксономий на русском языке с помощью больших языковых моделей'' методика TaxoLLaMA перенесена на русскоязычный материал RuWordNet. Авторы показывают, что дообученные LLM способны предсказывать гиперонимы, а модели, адаптированные под русский язык, в среднем дают результат лучше мультиязычных аналогов \cite{rcai2025,taxollama,ruwordnet}.

В этой курсовой похожая постановка переносится в более узкую область: ракетно-космическую технику. Вместо общего RuWordNet используется специализированный JSONL-дамп тезауруса РКТ. В нём каждому понятию соответствуют идентификатор, русские и английские синонимы, а также отношения с другими понятиями. Главный вопрос работы можно сформулировать так: можно ли использовать API-доступные LLM как генераторы гиперонимов для доменной технической онтологии, если после ответа сопоставлять сгенерированные фразы с идентификаторами тезауруса.

\textbf{Цель работы}: разработать и экспериментально оценить пайплайн применения LLM для предложения гиперонимов к терминам ракетно-космической техники с последующим сопоставлением ответов модели с идентификаторами доменной онтологии.

\textbf{Задачи работы:}
\begin{enumerate}
  \item изучить подход к автоматическому пополнению таксономий с помощью LLM на материале RuWordNet;
  \item подготовить доменный эталон на основе JSONL-дампа тезауруса ракетно-космической техники;
  \item реализовать конвертацию текстовых ответов LLM в идентификаторы понятий доменной онтологии;
  \item провести эксперименты с несколькими LLM и несколькими вариантами предметно-ориентированных промптов;
  \item сравнить результаты по метрикам MAP@15 и MRR@15, включая мягкую оценку по предкам второго уровня;
  \item разобрать основные причины ошибок и ограничения выбранной методики.
\end{enumerate}

\section*{Логика выполнения работы}
\addcontentsline{toc}{section}{Логика выполнения работы}

Работа выполнялась в несколько этапов. Сначала я попробовал воспроизвести общую схему экспериментов из статьи RCAI-2025: LLM получает термин и возвращает список гиперонимов, затем ответ преобразуется в идентификаторы тезауруса и оценивается по MAP@15 и MRR@15. Этот этап был скорее технической разведкой, чем отдельным экспериментом. Терминов было мало, полноценного доменного референса для РКТ ещё не было. Зато на нём удалось проверить формат ответов, разделитель между гиперонимами, конвертацию текстовых фраз в идентификаторы и подсчёт метрик. По смыслу этот шаг продолжает предыдущую курсовую работу, где проверялись методы машинного обучения и эмбеддинги для пополнения онтологий в области РКТ \cite{liss2025kurs}.

Потом был получен специализированный JSONL-дамп тезауруса ракетно-космической тематики \texttt{concepts.lii\_thes.kosmos.json}. После этого ручная подготовка эталона стала не нужна: референс строился автоматически из отношений \texttt{ВЫШЕ}. На этом этапе я собрал тестовую выборку из 50 терминов и реализовал маппинг ответа LLM в \texttt{conceptid} доменной онтологии.

Дальше к РКТ-выборке были применены те же типы промптов, что использовались в пилотной серии: \texttt{zero}, \texttt{cot} и \texttt{ontology}. Этот прогон ближе всего к идее воспроизведения RCAI-2025, только вместо общего RuWordNet используется специализированный тезаурус.

Отдельно были проверены промпты, явно настроенные на ракетно-космическую предметную область: \texttt{rkt\_basic}, \texttt{rkt\_few\_shot} и \texttt{rkt\_role}. Здесь я хотел проверить простую гипотезу: помогает ли модели указание предметной области и несколько коротких примеров из тезауруса.

\chapter{Постановка задачи}

Пусть задана онтология предметной области
\[
O = (C, R),
\]
где $C$ обозначает множество понятий, а $R$ обозначает множество отношений между ними. В этой работе рассматривается отношение \texttt{ВЫШЕ}. Оно задаёт переход от понятия к его гиперониму. Для каждого термина $t_i$ нужно получить упорядоченный список кандидатов:
\[
g(t_i) = \langle c_{i1}, c_{i2}, \ldots, c_{ik} \rangle,
\]
где $c_{ij} \in C$ — понятия онтологии, которые модель считает гиперонимами термина $t_i$. Такая формализация соответствует постановке задачи пополнения таксономий в работах по RuWordNet и TaxoLLaMA \cite{nikishina2022,rcai2025,taxollama}.

LLM в этой задаче не возвращает идентификаторы онтологии. Модель пишет обычные текстовые фразы: ``космический аппарат'', ``искусственный спутник Земли'', ``средство связи''. Поэтому пайплайн делится на два шага:
\begin{enumerate}
  \item генерация списка гиперонимов на естественном языке;
  \item сопоставление каждой сгенерированной фразы с \texttt{conceptid} в доменном тезаурусе.
\end{enumerate}

Итоговое качество зависит не только от языковой модели. Не меньше влияет нормализация ответа. Даже удачный текстовый кандидат может потеряться, если формулировки модели нет в синонимах тезауруса или она записана в ресурсе иначе.

\chapter{Обзор существующих подходов}

\section{Шаблонные методы извлечения гиперонимов}

Один из классических способов извлечения гиперонимии опирается на лексико-синтаксические шаблоны. Hearst предложила искать в корпусах конструкции вида ``$Y$, такие как $X$'' или ``$X$ и другие $Y$''. По таким конструкциям можно извлекать пары гипоним-гипероним из обычного текста \cite{hearst1992}. Если в корпусе встретилась фраза ``космические аппараты, такие как спутники'', из неё получается связь ``спутник'' $\rightarrow$ ``космический аппарат''.

Шаблоны хороши тем, что по ним сразу видно, откуда взялась связь. Но работают они только там, где родовое отношение выражено явно. В текстах о ракетно-космической технике это встречается не всегда. Чаще связь прячется в определении, описании назначения, перечислении характеристик или в аббревиатуре. Техническое описание может подробно раскрывать функцию разгонного блока, ни разу не используя фразу ``разгонный блок является...''. Поэтому шаблоны можно использовать как источник кандидатов, но строить на них всю доменную таксономию рискованно.

\section{RuThes, RuWordNet и пополнение таксономий}

RuWordNet представляет собой русский лексико-семантический ресурс, построенный на базе тезаурусной традиции RuThes \cite{loukachevitch2002,ruwordnet}. В таких ресурсах центральной единицей служит понятие или синсет. Ему соответствует набор лексических реализаций, а между понятиями задаются семантические отношения, включая гиперонимию. Эта идея близка к WordNet, где лексические значения организованы в синсеты и связаны семантическими отношениями \cite{miller1990}.

Лукашевич подчёркивает, что тезаурус должен хранить не только слова. В нём важны связи между понятиями: синонимические ряды, иерархические отношения и ассоциации \cite{louk2010}. Для пополнения ресурса это особенно существенно. Новый термин нужно встроить в правильное место графа, иначе он останется просто строкой в словаре.

В работах по пополнению RuWordNet задачу часто формулируют как встраивание новых слов или синсетов в уже существующую таксономию. В диахронической постановке ранняя версия ресурса используется как база, а понятия, появившиеся в более поздней версии, становятся тестовой выборкой \cite{nikishina2020,nikishina2022}. Такая схема даёт измеримую оценку: можно проверить, насколько хорошо метод предсказывает родительские узлы для новых терминов.

В этой курсовой сохраняется та же логика. Есть существующий граф понятий, и для выбранных терминов нужно предсказать родовые связи. Отличаются данные и технический контур. Вместо общего RuWordNet используется специализированный тезаурус РКТ, а вместо локально дообучаемых моделей — API-доступные LLM и промптинг.

\section{Векторные модели и графовые признаки}

До широкого применения LLM задачу пополнения таксономий часто решали через векторные представления слов и графовые признаки. Word2Vec и FastText позволяют измерять семантическую близость слов по их контекстам \cite{mikolov2013,bojanowski2017}. Для русского языка FastText удобен тем, что учитывает символьные n-граммы и лучше работает с редкими или морфологически сложными словами.

Более сильные методы используют не только текстовые эмбеддинги. В ход идут признаки из графа таксономии: глубина узла, число потомков, графовые представления, мета-эмбеддинги \cite{nikishina2022}. Для локального обучения и воспроизводимой оценки это хороший набор инструментов. В учебном проекте цена такого подхода довольно высокая: нужно готовить признаки, обучать модель и иметь полный доступ к ресурсу.

\section{Большие языковые модели и TaxoLLaMA}

TaxoLLaMA использует LLM для генерации гиперонимов и пополнения таксономических графов \cite{taxollama}. Замысел понятный: языковая модель уже содержит большой объём лексико-семантических знаний, а дообучение на парах гипоним-гипероним направляет эти знания на конкретную задачу. Эта линия работ связана с decoder-only моделями, включая семейство LLaMA \cite{touvron2023}.

Статья RCAI-2025 переносит эту методику на русский язык и RuWordNet \cite{rcai2025}. Авторы сравнивают мультиязычные модели и модели, адаптированные под русский язык, включая направление RuAdapt \cite{ruadapt2024}. После дообучения отдельные русскоязычные модели превосходят предыдущие векторные подходы по MAP. В статье также описана постобработка ответа модели: текстовые кандидаты нужно сопоставить с синсетами, иначе автоматическая оценка невозможна.

В этой курсовой LLM не дообучаются. Я рассматриваю более прикладной сценарий: API-модель генерирует кандидатов, а локальный код сопоставляет их с доменной онтологией РКТ. Для учебного и инженерного проекта это реалистичнее. Не нужно обучать собственную LLM, но остаются промпты, постобработка и строгая оценка.

\section{Место этой работы}

Эта работа находится между исследованием LLM для общего RuWordNet и прикладной задачей сопровождения специализированного тезауруса. По сравнению с RCAI-2025 меняются три компонента:
\begin{itemize}
  \item данные: вместо общего RuWordNet используется РКТ-тезаурус в JSONL;
  \item модели: вместо дообученных локальных моделей используются API-модели;
  \item главный технический риск: вместо обучения модели основной проблемой становится сопоставление свободного ответа LLM с \texttt{conceptid}.
\end{itemize}

Из-за этого в работе отдельно рассматриваются не только метрики качества. Важны маппинг, debug-файлы, soft-оценка и ошибки, которые появляются из-за расхождения между естественной формулировкой модели и внутренней структурой тезауруса.

\chapter{Данные}

\section{Доменный тезаурус РКТ}

Основной источник данных находится в файле \texttt{concepts.lii\_thes.kosmos.json}. Формат - JSONL: каждая строка содержит одно понятие тезауруса. В дампе 661 понятие. По структуре ресурс близок к тезаурусам семейства RuThes/RuWordNet: понятие задаётся идентификатором, набором лексических входов и семантическими отношениями с другими понятиями \cite{louk2010,loukachevitch2002,ruwordnet}.

Каждый объект включает поля:
\begin{itemize}
  \item \texttt{conceptid}: числовой идентификатор понятия;
  \item \texttt{conceptstr}: основное русское название понятия;
  \item \texttt{synonyms}: русские синонимы и лемматизированные формы;
  \item \texttt{e\_synonyms}: английские синонимы;
  \item \texttt{relats}: отношения с другими понятиями.
\end{itemize}

Запись читается следующим образом. Если объект имеет \texttt{conceptstr = "АЭРОКОСМИЧЕСКАЯ РАЗВЕДКА"}, это основное имя понятия. В массиве \texttt{synonyms} лежат русские варианты записи, включая исходную и лемматизированную формы. В массиве \texttt{relats} перечислены связи с другими понятиями. Элементы с \texttt{relationstr = "ВЫШЕ"} рассматриваются как родительские понятия для данной записи.

В эксперименте нужно всё время различать две вещи:
\begin{itemize}
  \item \textbf{текстовую форму}: то, что генерирует LLM, например ``космический аппарат'';
  \item \textbf{идентификатор понятия}: то, что требуется для формальной оценки, например \texttt{1023}.
\end{itemize}

Ответ модели сам по себе ещё не является результатом. Его нужно нормализовать, сопоставить с тезаурусом и превратить в список \texttt{conceptid}. Только после этого можно считать MAP@15 и MRR@15.

Для построения эталона используются только отношения, у которых \texttt{relationstr == "ВЫШЕ"}. Например, для понятия \texttt{АЭРОКОСМИЧЕСКАЯ РАЗВЕДКА} в дампе указаны родители \texttt{РАЗВЕДЫВАТЕЛЬНАЯ ДЕЯТЕЛЬНОСТЬ} и \texttt{КОСМИЧЕСКАЯ СЪЕМКА}.

\section{Построение эталонного референса}

Эталон строится автоматически из JSONL-дампа. Для каждого понятия $c_i$ извлекается список его родителей:
\[
Parents(c_i) = \{c_j \mid (c_i, \texttt{ВЫШЕ}, c_j) \in R\}.
\]

После этого формируется TSV-файл вида:
\begin{verbatim}
АЭРОКОСМИЧЕСКАЯ РАЗВЕДКА    ["2602", "4083"]
\end{verbatim}

Первый столбец содержит имя целевого термина. Второй столбец содержит JSON-список идентификаторов правильных родительских понятий. Такой формат выбран для совместимости со скриптом оценки \texttt{eval\_true.py}.

Автоматическое построение референса снимает сразу две проблемы. Во-первых, не приходится вручную выбирать гиперонимы для каждого термина. Во-вторых, эталон сразу получается в том формате, в котором проводится оценка. Ограничение при этом остаётся: если в тезаурусе у понятия указан только один непосредственный родитель, любой другой семантически близкий родитель в direct-оценке будет ошибкой. Именно поэтому в работе дополнительно используется soft-метрика.

\section{Тестовая выборка}

Для первичного эксперимента была сформирована тестовая выборка из 50 понятий, у которых есть хотя бы один родитель по отношению \texttt{ВЫШЕ}. Для этих понятий был построен TSV-референс в формате:
\begin{verbatim}
ТЕРМИН    ["parent_conceptid_1", "parent_conceptid_2", ...]
\end{verbatim}

Полный JSONL-дамп содержит 661 понятие. Я взял подвыборку из 50 терминов: её ещё можно разобрать вручную, но её уже достаточно, чтобы сравнить несколько моделей и промптов.

Размер выборки связан не только с удобством анализа. API-модель вызывается отдельно для каждого термина и каждого варианта промпта. Переход с 50 до 600 терминов увеличил бы число запросов более чем в десять раз. В рамках курсовой я выбрал другой баланс: несколько моделей и несколько промптов на управляемой выборке вместо одного большого прогона по всему тезаурусу.

В выборку вошли общие понятия и более конкретные названия объектов: типы летательных аппаратов, космические системы, спутники, элементы ракетно-космической техники и названия конкретных моделей. Для общих понятий модель должна выбрать правильный уровень абстракции. Для конкретных объектов трудность другая: модель не должна уходить в перечисление похожих объектов вместо родового класса.

\chapter{Метод}

\section{Общая схема и программный контур}

Перед основной серией была проведена пилотная ручная проверка схемы на малой выборке: 5 случайных слов и 5 вручную выбранных референсов из имеющихся файлов. Это не был полноценный эксперимент. Мне нужно было проверить формат ответа LLM, разделитель между гиперонимами, первичное сопоставление текстовых ответов с понятиями тезауруса и подсчёт MAP@15/MRR@15. После этого пайплайн был перенесён на специализированный JSONL-дамп РКТ-тезауруса, где референс строился автоматически по отношению \texttt{ВЫШЕ}.

Методическая часть работы строится вокруг одного процесса. Берётся термин из РКТ-тезауруса, LLM предлагает список возможных гиперонимов, текстовые ответы переводятся в идентификаторы понятий, после чего результат сравнивается с эталоном. Схема близка к постановке RCAI-2025: модель предсказывает гиперонимы, а затем предсказания приводятся к единицам RuWordNet \cite{rcai2025}. Здесь меняются ресурс и способ запуска. Вместо RuWordNet используется доменный JSONL-тезаурус РКТ, а вместо локального дообучения — API-модели.

Рабочий контур я собрал из нескольких скриптов. \texttt{build\_rkt\_reference.py} читает JSONL-дамп, извлекает связи \texttt{ВЫШЕ} и создаёт TSV-референс. \texttt{sample\_rkt\_dev\_terms.py} формирует тестовую выборку из понятий, у которых есть хотя бы один родитель. Запуск моделей выполняют \texttt{run\_llm\_batch.py} и \texttt{run\_llm\_batch\_rkt.py}. На выходе получается текстовый файл: в каждой строке стоит исходный термин и список гиперонимов, разделённых точкой с запятой.

После генерации выполняется постобработка. Скрипт \texttt{llm\_txt\_to\_predict\_rkt\_tsv.py} переводит текстовые гиперонимы в \texttt{conceptid}. Затем \texttt{eval\_true.py} считает direct MAP@15 и MRR@15, а \texttt{eval\_rkt\_soft.py} дополнительно считает soft-метрики с учётом предков на расстоянии не более двух шагов. Такое разделение полезно практически. Если итоговая метрика низкая, можно проверить, где произошла потеря: модель не предложила правильный гипероним, фраза не нашлась в индексе, нормализация сработала неверно или правильный кандидат оказался слишком низко в списке.

При построении индекса используется не только \texttt{conceptstr}, но и лексический слой понятия: \texttt{synonyms[*].textentrystr}, \texttt{synonyms[*].lementrystr}, при необходимости \texttt{e\_synonyms}. Отдельно индексируются \texttt{relats[*].conceptstr}. Этот пункт я добавил уже после первых прогонов. Оказалось, что часть родительских понятий встречается в JSONL только внутри отношений, а не как самостоятельная строка-концепт. Если такие строки не добавить в индекс, модель может назвать правильный родительский термин, но система не найдёт для него идентификатор.

\section{Маппинг текстовых ответов в \texttt{conceptid}}

LLM возвращает обычный текст. Например:
\begin{verbatim}
СПУТНИК СВЯЗИ    искусственный спутник Земли; космический аппарат; ретранслятор
\end{verbatim}
Человеку такая строка понятна. Для автоматической оценки она непригодна, потому что эталон хранит идентификаторы понятий. Каждый кандидат приходится нормализовать: привести к нижнему регистру, убрать лишние пробелы, унифицировать кавычки и дефисы. Для русских слов строится лемматизированная форма. Затем кандидат ищется в индексе тезауруса.

Этот шаг я недооценил на старте. Казалось, что строгого лексического сопоставления будет достаточно: модель назвала термин, индекс его нашёл, дальше считаются метрики. Первая же серия показала обратное. Часть нормальных по смыслу ответов просто терялась, потому что в тезаурусе такая фраза записана иначе или вообще отсутствует среди синонимов.

Если одна фраза соответствует нескольким \texttt{conceptid}, система сохраняет первые найденные варианты, но не меняет порядок, заданный моделью. В выходном TSV используются колонки \texttt{word}, \texttt{cand}, \texttt{cand\_name}, \texttt{prob}. Поле \texttt{prob} не является вероятностью модели. Это ранговый вес: первый найденный кандидат получает 1, второй получает $1/2$, третий $1/3$ и так далее. Такой вес нужен только для совместимости с форматом ранжированных предсказаний.

Схематически этап выглядит так:
\begin{verbatim}
for term in input_terms:
    answer = LLM(term, prompt)
    phrases = split(answer, ";")
    predictions = []
    for phrase in phrases:
        key = normalize_and_lemmatize(phrase)
        ids = term_index.get(key)
        append ids to predictions preserving rank
    write predictions to TSV
\end{verbatim}

Маппинг оказался отдельной исследовательской проблемой. Если модель пишет ``космический аппарат'', нужно определить, какому именно понятию тезауруса соответствует эта фраза. В ресурсе могут быть близкие варианты: ``космический аппарат'', ``автоматический космический аппарат'', ``пилотируемый космический корабль'', ``искусственный спутник Земли''. В технической области такие различия меняют место понятия в иерархии.

Показательный пример даёт термин \texttt{СПУТНИК СВЯЗИ}. Модель может предложить ``спутник-ретранслятор'' или ``космический ретранслятор''. По смыслу это нормальные варианты. Но если такой записи нет среди \texttt{conceptstr}, \texttt{synonyms} и \texttt{relats}, строгий маппинг не присвоит ей \texttt{conceptid}. В метрике это выглядит как ошибка, хотя содержательно модель могла попасть в нужную область. Именно для таких случаев в работе сохраняются debug-файлы: они показывают, какие фразы сопоставились с тезаурусом, а какие потерялись.

\section{Промпты, модели и протокол эксперимента}

Основная серия экспериментов построена по аналогии с RCAI-2025: проверяются разные способы инструкции модели без дообучения \cite{rcai2025}. Промпт \texttt{zero} просит прямо перечислить гиперонимы. Промпт \texttt{cot} просит модель сначала внутренне определить наиболее вероятные родовые понятия, но в ответ вывести только итоговый список. Промпт \texttt{ontology} ориентирует модель на краткие словарные наименования в стиле тезауруса. Все варианты требуют один формат ответа: гиперонимы через \texttt{;}, без комментариев.

Дополнительно были проверены РКТ-ориентированные инструкции. \texttt{rkt\_basic} прямо сообщает предметную область и запрещает слишком бытовые гиперонимы. \texttt{rkt\_few\_shot} даёт несколько коротких примеров из РКТ-тезауруса. \texttt{rkt\_role} задаёт роль составителя академического тезауруса. Эти промпты проверяют, помогает ли модели не просто знать общий смысл термина, а говорить в стиле доменного ресурса.

В эксперименте участвовали API-модели \texttt{gpt-5.4}, \texttt{claude-sonnet-4.6}, \texttt{claude-opus-4.7}, \texttt{deepseek-v4-pro}, \texttt{glm-5.1} и \texttt{kimi-k2.6}. Названия используются как API-идентификаторы конкретных прогонов. В основной таблице приведены 16 конфигураций. Каждый прогон обрабатывает 50 терминов, а каждый термин отправляется в отдельный API-запрос. С учётом моделей и промптов общий объём составил около полутора тысяч запросов API.

Для каждой пары ``модель/промпт'' выполнялась одна и та же процедура. Список терминов передавался в пакетный скрипт, скрипт формировал запросы к API, ответы сохранялись в текстовый файл, затем текст преобразовывался в TSV с \texttt{conceptid}. После этого считались direct- и soft-метрики.

Основные параметры генерации(!!!): температура 0.0 для воспроизводимости, максимальное число выходных токенов 400. Параметр \texttt{top\_p} явно не задавался; использовалось значение по умолчанию у API-провайдера. Нулевая температура нужна для повторяемости: при сравнении нескольких моделей и промптов важно, чтобы повторный запуск той же конфигурации давал тот же ответ. Иначе анализ ошибок быстро превращается в разбор случайных колебаний.

Модели не получали полный список понятий тезауруса на вход. Они свободно генерировали гиперонимы, а сопоставление с JSONL выполнялось уже после ответа. Это более жёсткий вариант, чем выбор из готового списка. Модель может написать удачную по смыслу формулировку, но ответ не будет зачтён, если такой строки нет в индексе.

\chapter{Метрики}

В работе используются две основные метрики: MAP@15 и MRR@15. В задачах пополнения таксономий они удобны, потому что оценивают не только наличие правильного гиперонима в ответе, но и его позицию в ранжированном списке кандидатов \cite{rcai2025,nikishina2022}.

Запись @15 означает, что оценка проводится только по первым 15 кандидатам в списке. Значение $K = 15$ выбрано по аналогии с работой RCAI-2025, где $K = 15$ задаёт максимальное число предсказаний при вычислении MAP@K и MRR@K. В этой работе сохраняется та же глубина ранжирования, чтобы результаты были методологически сопоставимы с исходной постановкой. Для практической работы с онтологией top-15 тоже выглядит разумно: эксперт получает достаточно широкий список, но выдача ещё не превращается в длинный набор почти случайных терминов.

Пусть для одного термина модель вернула ранжированный список кандидатов длиной не более $K$. Для позиции $k$ вводится индикатор
\[
I[y_k = 1],
\]
который равен 1, если кандидат на позиции $k$ является правильным гиперонимом, и равен 0 в противном случае.

Точность на первых $k$ позициях вычисляется так:
\[
P@k = \frac{1}{k}\sum_{j=1}^{k} I[y_j = 1].
\]
Эта величина показывает, какая доля кандидатов среди первых $k$ элементов списка оказалась правильной.

Для одного термина Average Precision вычисляется так:
\[
AP@K = \frac{1}{R}\sum_{k=1}^{K} P@k \cdot I[y_k = 1],
\]
где $R$ обозначает число правильных гиперонимов в эталоне. Вклад в сумму дают только те позиции, на которых модель действительно попала в правильный гипероним.

MAP@K является средним значением AP@K по всем терминам:
\[
MAP@K = \frac{1}{N}\sum_{i=1}^{N} AP@K_i.
\]

MRR@K оценивает позицию первого правильного ответа:
\[
MRR@K = \frac{1}{N}\sum_{i=1}^{N} \frac{1}{rank_i},
\]
где $rank_i$ обозначает позицию первого правильного кандидата для $i$-го термина. Если правильного кандидата в первых $K$ позициях нет, вклад этого термина в MRR считается равным 0.

Для анализа ошибок дополнительно используется hit@15. Эта метрика равна 1 для термина, если хотя бы один правильный гипероним попал в первые 15 кандидатов, и равна 0 иначе. Среднее значение hit@15 по выборке показывает долю терминов, для которых система нашла хотя бы один правильный ответ. В direct hit@15 учитываются только непосредственные родители, а в soft hit@15 также учитываются предки второго уровня.

Кроме прямой оценки используется \textbf{soft-оценка}. В ней релевантными считаются не только непосредственные родители понятия, но и его предки на расстоянии не более двух шагов по отношению \texttt{ВЫШЕ}. Похожая идея расширения эталона до гиперонимов второго порядка используется в диахронических наборах RuWordNet \cite{rcai2025}.

Direct-оценка проверяет, угадала ли модель именно тот родительский concept, который указан в тезаурусе как непосредственный. Soft-оценка смотрит мягче: попала ли модель хотя бы в близкую область иерархии. Для пополнения онтологии нужны обе оценки. Первая даёт строгое сравнение с эталоном, вторая помогает отделить грубые ошибки от случаев, где модель выбрала соседний уровень обобщения.

Soft MAP и soft MRR могут вести себя по-разному. При расширении эталона до предков второго уровня увеличивается число релевантных ответов $R$. В формуле AP@K сумма попаданий делится на $R$, поэтому soft MAP не обязана быть выше direct MAP. Soft MRR часто растёт, если хотя бы один близкий предок появляется в верхней части списка.

\chapter{Экспериментальные результаты}

{\scriptsize
\setlength{\tabcolsep}{2pt}
\begin{longtable}{p{29mm}p{18mm}p{14mm}rrrrr}
\caption{Результаты LLM на РКТ-выборке из 50 терминов\protect\footnotemark}\\
\toprule
Модель & Семейство & Промпт & Direct MAP & Direct MRR & Soft MAP & Soft MRR & Ср. канд. \\
\midrule
\endfirsthead
\toprule
Модель & Семейство & Промпт & Direct MAP & Direct MRR & Soft MAP & Soft MRR & Ср. канд. \\
\midrule
\endhead
gpt-5.4 & OpenAI & zero & 0.287 & 0.363 & 0.251 & 0.517 & 3.14 \\
gpt-5.4 & OpenAI & cot & 0.380 & 0.467 & 0.371 & 0.628 & 4.94 \\
gpt-5.4 & OpenAI & ontology & 0.214 & 0.256 & 0.252 & 0.381 & 4.24 \\
claude-sonnet-4.6 & Anthropic & zero & 0.260 & 0.320 & 0.183 & 0.410 & 2.58 \\
claude-sonnet-4.6 & Anthropic & cot & 0.302 & 0.391 & 0.251 & 0.489 & 4.26 \\
claude-sonnet-4.6 & Anthropic & ontology & 0.215 & 0.266 & 0.220 & 0.405 & 2.82 \\
claude-opus-4.7 & Anthropic & zero & 0.368 & 0.467 & 0.309 & 0.620 & н/д \\
claude-opus-4.7 & Anthropic & cot & 0.409 & 0.520 & 0.359 & 0.652 & 5.68 \\
claude-opus-4.7 & Anthropic & ontology & 0.237 & 0.291 & 0.247 & 0.407 & 3.76 \\
deepseek-v4-pro & DeepSeek & zero & 0.366 & 0.462 & 0.316 & 0.623 & 3.60 \\
deepseek-v4-pro & DeepSeek & cot & \textbf{0.422} & \textbf{0.550} & 0.379 & \textbf{0.675} & 5.56 \\
deepseek-v4-pro & DeepSeek & ontology & 0.198 & 0.246 & 0.225 & 0.344 & 4.26 \\
glm-5.1 & Zhipu/GLM & zero & 0.373 & 0.460 & 0.326 & 0.608 & 3.34 \\
glm-5.1 & Zhipu/GLM & cot & 0.379 & 0.516 & \textbf{0.382} & 0.646 & 5.10 \\
kimi-k2.6 & Moonshot/Kimi & zero & 0.353 & 0.466 & 0.304 & 0.587 & 3.30 \\
kimi-k2.6 & Moonshot/Kimi & cot & 0.320 & 0.416 & 0.307 & 0.557 & 4.60 \\
\bottomrule
\end{longtable}
}
\footnotetext{Для моделей \texttt{glm-5.1} и \texttt{kimi-k2.6} промпт \texttt{ontology} не приведён. При попытке прогона API-провайдер \texttt{ohmylama} массово возвращал таймауты и ошибки 502. После повторных запусков эти конфигурации всё равно не завершились, поэтому в таблицу не добавлены псевдо-результаты.}


Лучший результат по строгим direct-метрикам дала конфигурация \texttt{deepseek-v4-pro + cot}: MAP@15 = 0.422 и MRR@15 = 0.550. В этой связке непосредственный родительский концепт часто оказывается в верхней части списка кандидатов.

По Soft MAP@15 лучший результат у \texttt{glm-5.1 + cot}: 0.382. По Soft MRR@15 снова лидирует \texttt{deepseek-v4-pro + cot}: 0.675. Такая разница возникает из-за того, что модели нередко предлагают не прямого родителя, а более общий предок, расположенный рядом с целевым понятием в таксономии.

\section{Интерпретация результатов}

Для первичного эксперимента результаты получились содержательными. Без дообучения LLM и без полного списка понятий на входе лучшие конфигурации достигают MAP@15 выше 0.42 на строгой оценке. Модель не просто выдаёт общие слова. Во многих случаях её ответ удаётся сопоставить с идентификатором доменного тезауруса.

Самым устойчивым оказался промпт \texttt{cot}. Почти для всех моделей он улучшает direct- и soft-метрики по сравнению с zero-shot. Исключение составляет \texttt{kimi-k2.6}: у этой модели zero-shot сильнее по direct MAP и direct MRR. Промпт \texttt{ontology} в этой серии чаще ухудшал результат. Вероятная причина в том, что слишком жёсткая ориентация на словарные формулировки не совпадала со структурой конкретного РКТ-тезауруса и приводила к потере релевантных кандидатов.

Средние значения по типам промптов:
\begin{itemize}
  \item \texttt{zero}: direct MAP@15 = 0.335, direct MRR@15 = 0.423;
  \item \texttt{cot}: direct MAP@15 = 0.369, direct MRR@15 = 0.477;
  \item \texttt{ontology}: direct MAP@15 = 0.216, direct MRR@15 = 0.265.
\end{itemize}

Инструкция \texttt{cot} не заставляет модель показывать рассуждения в ответе. Наоборот, промпт требует вывести только итоговый список. Но просьба сначала внутренне определить наиболее вероятные гиперонимы помогает лучше упорядочить кандидатов. Для MRR это заметно сильнее всего: метрика чувствительна к позиции первого правильного ответа.

Сравнивать семейства моделей нужно осторожно. Для \texttt{Zhipu/GLM} и \texttt{Moonshot/Kimi} было меньше конфигураций, чем для OpenAI и Anthropic. Поэтому корректнее смотреть на отдельные связки ``модель/промпт''. \texttt{deepseek-v4-pro} дал лучший одиночный direct-результат, а \texttt{glm-5.1} лучший soft MAP.

Soft-MRR обычно выше direct-MRR. Модель часто находит родовое понятие рядом с правильным ответом в иерархии, но не совпадает с непосредственным родителем. Soft-MAP иногда ниже direct-MAP, потому что при расширении эталона увеличивается число релевантных элементов $R$, а AP делит сумму попаданий на $R$. В этой работе soft-MRR особенно полезна как диагностическая метрика: она показывает, насколько высоко в списке появляется хотя бы один близкий предок.

\section{Дополнительная серия с РКТ-промптами}

После основной серии были проведены дополнительные эксперименты с промптами, ориентированными на терминологию ракетно-космической техники. В этой серии участвовали три модели: \texttt{deepseek-v4-pro}, \texttt{claude-opus-4.7} и \texttt{gpt-5.4}. Проверялись три варианта инструкции:
\begin{itemize}
  \item \texttt{rkt\_basic}: прямое указание предметной области и запрет на слишком бытовые гиперонимы;
  \item \texttt{rkt\_few\_shot}: несколько примеров из РКТ-тезауруса;
  \item \texttt{rkt\_role}: роль составителя академического тезауруса.
\end{itemize}

{\scriptsize
\setlength{\tabcolsep}{3pt}
\begin{longtable}{p{20mm}p{26mm}rrrrr}
\caption{Дополнительные результаты с РКТ-ориентированными промптами}\\
\toprule
Модель & Промпт & Direct MAP & Direct MRR & Soft MAP & Soft MRR & Ср. канд. \\
\midrule
\endfirsthead
\toprule
Модель & Промпт & Direct MAP & Direct MRR & Soft MAP & Soft MRR & Ср. канд. \\
\midrule
\endhead
dsv4 & rkt\_basic & 0.274 & 0.363 & 0.212 & 0.431 & 2.92 \\
dsv4 & rkt\_few\_shot & \textbf{0.417} & \textbf{0.572} & 0.351 & 0.635 & 4.60 \\
dsv4 & rkt\_role & 0.345 & 0.480 & 0.262 & 0.540 & 2.58 \\
opus47 & rkt\_basic & 0.279 & 0.380 & 0.255 & 0.500 & 2.30 \\
opus47 & rkt\_few\_shot & 0.410 & 0.550 & \textbf{0.366} & \textbf{0.677} & 4.10 \\
opus47 & rkt\_role & 0.165 & 0.227 & 0.157 & 0.323 & 2.72 \\
gpt54 & rkt\_basic & 0.279 & 0.365 & 0.235 & 0.495 & 2.66 \\
gpt54 & rkt\_few\_shot & 0.346 & 0.420 & 0.254 & 0.545 & 2.44 \\
gpt54 & rkt\_role & 0.360 & 0.500 & 0.295 & 0.640 & 2.26 \\
\bottomrule
\end{longtable}
}

РКТ-ориентированная серия показала, что примеры из предметной области могут улучшать качество. Самый показательный вариант — \texttt{rkt\_few\_shot}. Он дал лучшие результаты для \texttt{dsv4} и \texttt{opus47}. Для \texttt{gpt-5.4} сильнее по direct MRR и soft MRR оказался \texttt{rkt\_role}; роль составителя тезауруса лучше согласовала ответы этой модели со структурой онтологии.

Но доменная инструкция не выигрывает автоматически у общей серии \texttt{zero/cot/ontology}. Например, для \texttt{deepseek-v4-pro} общий \texttt{cot}-промпт дал direct MAP@15 = 0.422, а \texttt{rkt\_few\_shot} показал 0.417. Разница небольшая, но она важна: простое добавление предметной области не всегда сильнее хорошо сформулированного общего промпта. Качество зависит от сочетания модели, инструкции и последующего маппинга ответа в онтологию.

Отдельно бросается в глаза \texttt{opus47 + rkt\_role}. Я перепроверял этот прогон, потому что просадка с уровня около 0.41 до 0.165 выглядит слишком резкой. Но по сохранённым выдачам видно, что модель действительно стала хуже попадать в прямых родителей: инструкция сузила ответы и испортила ранжирование.

Чтобы сравнить промпты в равных условиях, отдельно рассмотрены модели, у которых выполнены все шесть вариантов инструкции: \texttt{deepseek-v4-pro}, \texttt{gpt-5.4} и \texttt{claude-opus-4.7}. Для них direct MAP@15 по промптам приведён в таблице~\ref{tab:prompts_h2h}.

\begin{table}[H]
\centering
\caption{Direct MAP@15 по промптам для моделей, у которых выполнены все шесть вариантов}
\label{tab:prompts_h2h}
\begin{tabular}{lcccc}
\toprule
Промпт & \texttt{dsv4} & \texttt{gpt54} & \texttt{opus47} & Среднее \\
\midrule
\texttt{cot}            & 0.422 & 0.380 & 0.409 & 0.404 \\
\texttt{rkt\_few\_shot} & 0.417 & 0.346 & 0.410 & 0.391 \\
\texttt{zero}           & 0.366 & 0.287 & 0.368 & 0.340 \\
\texttt{rkt\_role}      & 0.345 & 0.360 & 0.165 & 0.290 \\
\texttt{rkt\_basic}     & 0.274 & 0.279 & 0.279 & 0.277 \\
\texttt{ontology}       & 0.198 & 0.214 & 0.237 & 0.216 \\
\bottomrule
\end{tabular}
\end{table}

В этой подвыборке порядок промптов довольно устойчив. Общий \texttt{cot} в среднем сильнее доменного \texttt{rkt\_few\_shot} на 0.013, причём для двух моделей из трёх \texttt{cot} выигрывает напрямую. На \texttt{deepseek-v4-pro} разница оказывается в третьем знаке. Доменные примеры полезны, но в текущей постановке они не дают системного преимущества над хорошей общей формулировкой.

\section{Сравнение основной и РКТ-ориентированной серий}

Две серии экспериментов отвечают на разные вопросы. Основная серия \texttt{zero/cot/ontology} проверяет, переносится ли общая методика предсказания гиперонимов из RCAI-2025 на РКТ-тезаурус. РКТ-ориентированная серия проверяет, помогает ли явно сообщать модели предметную область и давать примеры из неё.

По результатам видно, что доменная инструкция полезна, но не гарантирует улучшение. \texttt{dsv4 + rkt\_few\_shot} почти повторяет качество \texttt{deepseek-v4-pro + cot}. Для \texttt{opus47} доменный few-shot оказался сильным по soft-метрикам. При этом \texttt{rkt\_role} для той же модели резко ухудшил результат. Слишком жёсткая инструкция может сузить пространство ответов и испортить ранжирование.

Практический вывод здесь простой: при использовании LLM для пополнения онтологии нельзя ограничиваться одной формулировкой промпта. Даже близкие по смыслу инструкции дают разные списки кандидатов. Промпт в этом эксперименте оказался полноценным фактором, а не мелкой технической настройкой.

\chapter{Анализ ошибок}

Анализ ошибок проводился по debug-файлам, которые сохранялись после маппинга ответов LLM в \texttt{conceptid}. В этих файлах виден не только итоговый промах или попадание. Видно, какую фразу написала модель, удалось ли найти её в тезаурусе, какой идентификатор был выбран и на какой позиции он оказался. Такой разбор полезнее простой таблицы метрик: одинаковая ошибка в MAP может появиться по разным причинам.

Первая группа ошибок связана с потерями на этапе сопоставления. Модель может предложить фразу, которая по смыслу подходит, но отсутствует в индексе онтологии. Для термина \texttt{СПУТНИК СВЯЗИ} встречаются ответы ``спутник-ретранслятор'', ``орбитальное средство связи'', ``космический ретранслятор''. Для специалиста это естественные варианты, но они не всегда совпадают с \texttt{conceptstr} или синонимами в JSONL. Такая фраза не получает \texttt{conceptid} и не участвует в оценке. Метрика падает, хотя проблема возникла уже после генерации.

Вторая группа ошибок связана с уровнем обобщения. Иногда модель выбирает слишком общий гипероним: ``система'', ``объект'', ``техническое устройство''. Такие ответы могут быть разумными на уровне языка, но в тезаурусе они находятся слишком далеко от непосредственного родителя. Для конкретного беспилотного аппарата ответ ``летательный аппарат'' выглядит нормальным. Но если в тезаурусе прямым родителем является ``беспилотный летательный аппарат'' или более узкая категория, direct-оценка засчитает ошибку. Soft-оценка частично помогает, потому что учитывает близких предков.

Бывает и обратная ситуация: модель предлагает соседнее или дочернее понятие вместо родителя. Это особенно заметно на терминах с названиями конкретных летательных аппаратов, ракет и спутников. Модель узнаёт предметную область и начинает перечислять похожие классы техники, но не всегда соблюдает направление отношения. Для семантического поиска такой ответ ещё мог бы быть полезен. Для пополнения таксономии этого мало: нужно различать родителя, потомка и просто соседний термин из той же ветви.

Часть ошибок объясняется устройством самого тезауруса. Он может фиксировать такую структуру, которая не совпадает с первым языковым ожиданием человека. В работах по онтологиям и тезаурусам это описывается как риск смешения родовидовых отношений, ролей, отношений часть-целое и класса-экземпляра \cite{louk2010,ontotez2006}. В РКТ-тезаурусе это проявляется на технических объектах и названиях моделей. Правильный родитель может быть точным с точки зрения структуры ресурса, но не самым очевидным для свободной генерации.

Есть и лексическая причина. JSONL-дамп содержит большой набор синонимов, но он не покрывает все возможные формулировки. LLM может использовать современное обозначение, разговорный вариант, английскую аббревиатуру или смешанную запись. Для форм вроде \texttt{F-16}, \texttt{MQ-35A}, \texttt{J-10} обычная русская лемматизация почти не помогает. С многословными терминами проблема ещё заметнее: перестановка слов или замена одного компонента может сделать фразу невидимой для точного поиска.

Последний тип ошибок связан с ранжированием. Правильный \texttt{conceptid} иногда есть в списке, но стоит слишком низко. Для MRR это критично: первый правильный ответ на первой позиции даёт вклад 1, на пятой позиции только 0.2. Поэтому система должна генерировать не просто много кандидатов, а ставить наиболее вероятные родовые понятия в начало. По результатам видно, что увеличение среднего числа кандидатов само по себе качества не гарантирует.

\section{Количественная оценка достижимости}

Перечисленные группы ошибок объясняют причины промахов, но не показывают, какая доля выборки вообще достижима текущим пайплайном. Для такой оценки все 25 конфигураций , 16 в основной серии и 9 в РКТ-ориентированной. Были рассмотрены как ансамбль предсказателей на одной и той же выборке из 50 терминов. Для каждого термина введён бинарный признак: попал ли хотя бы один правильный родитель в top-15 у данной конфигурации. Затем подсчитано, сколько конфигураций решают каждый термин и сколько терминов решает каждое подмножество конфигураций.

Из 50 терминов 40 решает хотя бы одна конфигурация, а 10 не решает ни одна. Верхняя граница hit@15 на этой выборке составляет 0.80. Лучшая единичная конфигурация в эксперименте даёт hit@15 = 0.68. Значит, ансамблирование может закрыть зазор не более чем на 0.12. Оставшиеся 0.20 не закрываются ни одной из проверенных связок: правильный \texttt{conceptid} не появляется в top-15 ни у одной из 25 конфигураций.

Шесть из десяти нерешённых терминов остаются нерешёнными даже в режиме soft с предками на расстоянии не более двух шагов. Это \texttt{АВИАЦИОННЫЙ ПОЛЕТ}, \texttt{АО ``КАМОВ''}, \texttt{ВЕРТОЛЕТ МИ-8}, \texttt{ПОСАДКА ЛЕТАТЕЛЬНОГО АППАРАТА}, \texttt{РАКЕТА ``СТУГНА-П''} и \texttt{СПУТНИКОВЫЙ ТЕЛЕФОН}. Debug-файлы показывают однотипную картину. Модели предлагают разумные родовые формулировки: ``телефон'', ``средство связи'', ``посадка'', ``этап полёта'', ``маневр воздушного судна''. Индекс тезауруса их не разрешает. Например, для термина \texttt{СПУТНИКОВЫЙ ТЕЛЕФОН} конфигурация \texttt{deepseek-v4-pro + cot} вернула 15 кандидатов: ``телефон'', ``средство связи'', ``беспроводной телефон'', ``радиотелефон'' и подобные. Ни один не получил \texttt{conceptid}. Для \texttt{ПОСАДКА ЛЕТАТЕЛЬНОГО АППАРАТА} картина такая же: 15 кандидатов и 0 совпадений с индексом. Это ограничение текущего сопоставления свободного текста с идентификаторами. Простая замена модели его не снимает.

Распределение терминов по числу решающих их конфигураций получилось выраженно бимодальным (таблица~\ref{tab:term_difficulty}). На одном краю находятся 19 терминов, для которых правильный родитель попадает в top-15 у 21--25 конфигураций из 25; такие термины уверенно решает почти любая разумная связка модель-промпт. На другом краю находятся те же 10 терминов с нулевым попаданием. Середина распределения, где выбор конфигурации действительно влияет на результат, занимает меньшую часть выборки. На этой выборке перебор моделей и промптов помогает, но ресурс такого улучшения быстро заканчивается.

Согласованность конфигураций оценивалась попарной корреляцией Пирсона по вектору per-term AP@15. Среднее значение составило 0.627, медиана — 0.626. Минимум 0.259 даёт пара \texttt{claude-sonnet-4.6 + zero} и \texttt{claude-opus-4.7 + rkt\_role}, максимум 0.966 — пара \texttt{deepseek-v4-pro + zero} и \texttt{glm-5.1 + zero}. Конфигурации не идентичны, но и независимыми их назвать нельзя. В реальной системе, скорее всего, полезнее сочетать разные типы инструкций, чем много раз варьировать одну и ту же идею.

Сводные количественные ориентиры приведены в таблице~\ref{tab:upper_bound}. Там же добавлено значение oracle direct MAP@15: для каждого термина выбран конфиг с наибольшим AP@15, и эти значения усреднены по выборке. Это теоретический потолок ансамбля при идеальном переключении конфигураций по терминам. Лучший единичный direct hit@15 равен 0.68: этого значения достигают пять конфигураций (\texttt{deepseek-v4-pro + cot}, \texttt{deepseek-v4-pro + rkt\_few\_shot}, \texttt{gpt-5.4 + cot}, \texttt{claude-opus-4.7 + cot} и \texttt{claude-opus-4.7 + ontology}); каждая из них закрывает 34 термина из 50, а наибольший MAP@15 среди них у \texttt{deepseek-v4-pro + cot}.

\begin{table}[H]
\centering
\caption{Верхняя граница пайплайна на 50 РКТ-терминах при объединении 25 конфигураций}
\label{tab:upper_bound}
\begin{tabular}{p{115mm}r}
\toprule
 & Значение \\
\midrule
Лучший единичный direct MAP@15 (\texttt{deepseek-v4-pro + cot}) & 0.422 \\
Лучший единичный direct hit@15 & 0.68 \\
Oracle direct MAP@15 (выбор лучшей конфигурации для каждого термина) & 0.566 \\
Upper-bound direct hit@15 (хотя бы одна конфигурация) & 0.80 \\
Upper-bound soft hit@15 (depth $\leq$ 2) & 0.88 \\
Терминов без попадания у всех 25 конфигураций (direct) & 10/50 \\
Терминов без попадания у всех 25 конфигураций (soft) & 6/50 \\
Средняя попарная корреляция Пирсона per-term AP@15 & 0.627 \\
\bottomrule
\end{tabular}
\end{table}

Отдельно проверялся простой жадный отбор конфигураций по приросту объединённого hit@15. На первом шаге выбирается конфигурация с наибольшим hit@15, дальше к ней добавляются те конфигурации, которые добавляют к покрытию наибольшее число ранее нерешённых терминов. Три конфигурации — \texttt{deepseek-v4-pro + cot}, \texttt{deepseek-v4-pro + rkt\_basic} и \texttt{claude-sonnet-4.6 + cot} — дают объединённый hit@15 = 0.78 при общей верхней границе 0.80. Почти всё достижимое покрытие можно получить малым ансамблем. Дальнейшее добавление близких конфигураций даёт небольшой прирост.

\begin{table}[H]
\centering
\caption{Распределение терминов по числу конфигураций, у которых правильный родитель попадает в top-15 (direct)}
\label{tab:term_difficulty}
\begin{tabular}{cr}
\toprule
Число решающих конфигураций & Число терминов \\
\midrule
0      & 10 \\
1--5   &  3 \\
6--10  &  5 \\
11--15 &  5 \\
16--20 &  8 \\
21--25 & 19 \\
\bottomrule
\end{tabular}
\end{table}

Разбор ошибок приводит к довольно практическому выводу. Успех системы зависит не только от того, ``знает'' ли LLM термин. Нужны правильный родовой кандидат, совпадение с лексическим слоем тезауруса, устойчивая нормализация и адекватное ранжирование. Дальше стоит усиливать не только промпты и выбор моделей, но и постобработку: расширение синонимов, поиск по эмбеддингам, подсказки с кандидатами из тезауруса и LLM-переранжирование. Идея использовать LLM как переранжировщик уже обсуждается в работах по информационному извлечению \cite{ma2023reranker}.

\chapter{Сравнение с работой RCAI-2025}

Работа RCAI-2025 исследует дообучение LLM на парах гипоним-гипероним из RuWordNet. Там основной акцент сделан на переносе TaxoLLaMA на русский язык и сравнении мультиязычных и русскоязычных моделей \cite{rcai2025,taxollama}. Эта курсовая решает смежную, но более прикладную задачу.

\begin{itemize}
  \item В RCAI-2025 используется общий лексико-семантический ресурс RuWordNet. Здесь используется доменный тезаурус ракетно-космической техники.
  \item В RCAI-2025 модели дообучаются на большом наборе пар гипоним-гипероним. Здесь исследуются API-модели без дообучения, но с предметно-ориентированными промптами.
  \item В RCAI-2025 главный вызов связан с адаптацией LLM к русскоязычной таксономии. Здесь главный вызов связан с сопоставлением свободного ответа LLM с идентификаторами специализированной онтологии.
\end{itemize}

Абсолютные значения MAP/MRR напрямую сравнивать некорректно: различаются данные, модели, способ обучения и процедура постобработки. Методологическая связь при этом сохраняется. В обоих случаях задача сводится к автоматическому пополнению таксономии через предсказание гиперонимов.

Есть и различие в масштабе. RCAI-2025 рассматривает дообучение локальных моделей на большом наборе пар гипоним-гипероним. В этой работе использовались API-модели, а каждый термин обрабатывался отдельным запросом. Выборка из 50 терминов стала компромиссом: она достаточна для сравнения нескольких моделей и промптов, но остаётся управляемой по стоимости и времени.

\chapter{Ограничения}

Проведённый эксперимент имеет несколько ограничений:
\begin{itemize}
  \item тестовая выборка содержит 50 терминов, поэтому результаты нужно читать как первичный эксперимент, а не как финальную оценку метода;
  \item использовались API-модели без дообучения, поэтому влияние обучения на парах из доменной онтологии не проверялось;
  \item сопоставление текстовых гиперонимов с \texttt{conceptid} основано в основном на нормализации и лемматизации, без полноценного семантического поиска;
  \item в JSONL часть родительских понятий присутствует только в отношениях, что усложняет построение полного индекса синонимов;
  \item direct-метрика может занижать качество ответов, если модель предлагает близкий, но не непосредственный родительский concept.
\end{itemize}

Сэмплинг был ограниченным. Никакой стратификации по типам терминов в этой версии не делалось. Это тоже важно учитывать: 50 терминов позволяют увидеть общую картину, но не дают надёжно сравнивать, например, организации, процессы и конкретные модели техники как отдельные классы примеров.

\chapter{Перспективы развития}

Полученный пайплайн можно развивать как цельную систему пополнения онтологии, а не как набор отдельных скриптов. Сейчас модель генерирует гиперонимы свободно, а локальный код пытается сопоставить эти фразы с тезаурусом. Для первого эксперимента этого достаточно: быстро проверяется, применимы ли LLM к доменной задаче. Но главный запас качества лежит в связке ``генерация плюс поиск''. Чем лучше система ограничивает пространство возможных ответов, тем меньше потерь на маппинге.

Количественная оценка достижимости из главы ``Анализ ошибок'' даёт конкретные ориентиры для следующих шагов. Oracle direct MAP@15 = 0.566 при лучшем единичном результате 0.422 показывает, что ансамбль конфигураций имеет ограниченный, но измеримый запас. Жадный отбор по объединённому hit@15 показал, что три конфигурации (\texttt{deepseek-v4-pro + cot}, \texttt{deepseek-v4-pro + rkt\_basic} и \texttt{claude-sonnet-4.6 + cot}) уже дают объединённый hit@15 = 0.78 при верхней границе 0.80. Дальнейшее добавление близких по поведению конфигураций почти не расширяет покрытие. Поэтому я бы дальше делал ставку не на бесконечный перебор LLM, а на расширение лексического слоя тезауруса и семантический fallback при сопоставлении ответов с \texttt{conceptid}. Именно там сейчас лежит зазор в 0.20 hit@15, который приходится на структурно недостижимые термины.

Первое направление развития связано с расширением сопоставления фраз LLM с \texttt{conceptid}. В текущей версии используются нормализация, лемматизация и поиск по лексическому индексу. Следующий шаг, добавить семантический поиск: если точного совпадения нет, система должна искать ближайшие понятия тезауруса по эмбеддингам. Например, ответ ``спутник-ретранслятор'' можно сопоставлять не только по строке, но и по близости к понятиям ``спутник связи'', ``ретранслятор'', ``искусственный спутник Земли''. Такой подход близок к retrieval-augmented generation, где генерация комбинируется с поиском по внешнему ресурсу \cite{lewis2020rag}.

Второе направление связано с изменением роли LLM. Вместо свободной генерации можно сначала получить список кандидатов из тезауруса: по строковому поиску, эмбеддингам, графовой близости или частотным признакам. После этого LLM получает термин и уже готовый список \texttt{conceptid} с названиями понятий. Задача модели становится другой: не придумать гипероним с нуля, а ранжировать кандидатов. Такая схема устойчивее, потому что модель не выходит за пределы тезауруса и не создаёт фразы, которые потом невозможно оценить. В задачах с трудными кандидатами LLM уже показывают себя полезными переранжировщиками \cite{ma2023reranker}.

Изначально я хотел быстрее перейти к дообучению на доменном материале. После анализа ошибок это стало выглядеть преждевременным. Если маппинг нестабилен, улучшение модели трудно отделить от улучшения постобработки. Поэтому дообучение имеет смысл делать уже после того, как будет построен более надёжный референс и добавлен хотя бы простой семантический поиск. Тогда можно подготовить пары вида ``термин, гипероним'' из JSONL и обучить модель или компактный ранжировщик на этих примерах.

Четвёртое направление связано с расширением тестовой выборки. Сейчас используются 50 терминов. Этого достаточно для первичной проверки нескольких моделей и промптов, но мало для анализа по типам терминов. Следующая версия эксперимента может включать 100-150 понятий и стратификацию: конкретные объекты, классы техники, процессы, организации, абстрактные понятия. Тогда можно будет увидеть, где LLM ошибается чаще. Например, для конкретных моделей техники вероятны ошибки соседства, а для абстрактных понятий чаще возникает слишком общий уровень обобщения.

Наиболее практичной выглядит гибридная схема. Поиск по тезаурусу даёт ограниченный набор кандидатов, LLM ранжирует их с учётом значения термина, а soft-оценка показывает, попадает ли ответ хотя бы в правильную ветвь иерархии. Такая архитектура ближе к реальной задаче сопровождения онтологии: система не заменяет эксперта, а готовит для него короткий список проверяемых родительских понятий с понятными основаниями выбора.

\chapter*{Заключение}
\addcontentsline{toc}{chapter}{Заключение}

В работе реализован и протестирован пайплайн применения больших языковых моделей для пополнения доменной онтологии ракетно-космической техники. В отличие от экспериментов на RuWordNet, здесь используется специализированный JSONL-дамп тезауруса РКТ, из которого автоматически строится эталон по отношению \texttt{ВЫШЕ}.

В ходе работы были реализованы:
\begin{itemize}
  \item построение РКТ-референса из JSONL;
  \item формирование тестовой выборки из 50 терминов;
  \item запуск нескольких LLM через API;
  \item предметно-ориентированные промпты;
  \item конвертация текстовых ответов LLM в \texttt{conceptid};
  \item direct- и soft-оценка качества.
\end{itemize}

В основной серии \texttt{zero/cot/ontology} лучший direct-результат получен для конфигурации \texttt{deepseek-v4-pro + cot}: MAP@15 = 0.422, MRR@15 = 0.550. Лучший Soft MAP@15 получен для \texttt{glm-5.1 + cot}: 0.382, а лучший Soft MRR@15 получен для \texttt{deepseek-v4-pro + cot}: 0.675.

В дополнительной серии РКТ-ориентированных промптов лучший direct-результат дала конфигурация \texttt{deepseek-v4-pro + rkt\_few\_shot}: MAP@15 = 0.417, MRR@15 = 0.572. Лучший soft-результат получен для \texttt{claude-opus-4.7 + rkt\_few\_shot}: Soft MAP@15 = 0.366, Soft MRR@15 = 0.677.

Эксперимент показывает, что LLM можно использовать как генераторы кандидатов для пополнения специализированной онтологии. Но итоговое качество сильно зависит от сопоставления текстового ответа с идентификаторами онтологии. Самый перспективный следующий шаг связан с гибридной схемой: LLM генерирует или ранжирует кандидатов, а поиск по тезаурусу и корпусу РКТ ограничивает пространство возможных гиперонимов.

\addcontentsline{toc}{chapter}{Список литературы}
\begin{thebibliography}{99}

\bibitem{gruber1993}
Gruber T. R.
A translation approach to portable ontology specifications.
Knowledge Acquisition. 1993. Vol. 5(2). pp. 199-220.

\bibitem{louk2010}
Лукашевич Н.\,В.
Тезаурусы в задачах информационного поиска.
М., 2010. 396 с.

\bibitem{ontotez2006}
Соловьев В.\,Д., Добров Б.\,В., Иванов В.\,В., Лукашевич Н.\,В.
Онтологии и тезаурусы: учебное пособие.
Казань; Москва, 2006.

\bibitem{liss2025kurs}
Лисс М.\,А.
Методы машинного обучения для автоматического пополнения онтологий в специальной предметной области разработки ракетно-космической техники.
Курсовая работа. МГУ имени М.\,В.\,Ломоносова, Факультет космических исследований. Москва, 2025.

\bibitem{rcai2025}
Садковский Ф.\,А., Лукашевич Н.\,В., Гришин И.\,Ю.
Автоматическое пополнение таксономий на русском языке с помощью больших языковых моделей.
Конференция по искусственному интеллекту RCAI, 2025.

\bibitem{hearst1992}
Hearst M. A.
Automatic acquisition of hyponyms from large text corpora.
Proceedings of COLING, 1992.

\bibitem{miller1990}
Miller G. A., Beckwith R., Fellbaum C., Gross D., Miller K. J.
Introduction to WordNet: An on-line lexical database.
International Journal of Lexicography. 1990. Vol. 3(4). pp. 235-244.

\bibitem{ruwordnet}
Loukachevitch N. V. et al.
Creating Russian WordNet by conversion.
Computational Linguistics and Intellectual Technologies, 2016.

\bibitem{loukachevitch2002}
Loukachevitch N. V., Dobrov B. V.
Development and use of thesaurus of Russian language RuThes.
Proceedings of Workshop on WordNet Structures and Standardisation, 2002.

\bibitem{taxollama}
Moskvoretskii V., Nikishina I., Neminova E., Lobanova A., Panchenko A., Biemann C.
Large Language Models for Creation, Enrichment and Evaluation of Taxonomic Graphs.
Semantic Web. Vol.~17(1). pp.~1--30.
DOI: 10.1177/22104968251404186.

\bibitem{nikishina2022}
Nikishina I. et al.
Taxonomy enrichment with text and graph vector representations.
Semantic Web, 2022.

\bibitem{nikishina2020}
Nikishina I., Logacheva V., Panchenko A., Loukachevitch N.
RUSSE'2020: Findings of the First Taxonomy Enrichment Task for the Russian language.
Computational Linguistics and Intellectual Technologies, 2020.

\bibitem{mikolov2013}
Mikolov T., Chen K., Corrado G., Dean J.
Efficient estimation of word representations in vector space.
arXiv:1301.3781, 2013.

\bibitem{bojanowski2017}
Bojanowski P., Grave E., Joulin A., Mikolov T.
Enriching word vectors with subword information.
Transactions of the Association for Computational Linguistics. 2017. Vol. 5. pp. 135-146.

\bibitem{touvron2023}
Touvron H. et al.
LLaMA: Open and efficient foundation language models.
arXiv:2302.13971, 2023.

\bibitem{ruadapt2024}
Tikhomirov M. M., Chernyshev D. I.
Improving Large Language Model Russian adaptation with preliminary vocabulary optimization.
Lobachevskii Journal of Mathematics. 2024. Vol. 45(7). pp. 3211-3219.

\bibitem{ma2023reranker}
Ma Y., Cao Y., Hong Y., Sun A.
Large Language Model Is Not a Good Few-shot Information Extractor, but a Good Reranker for Hard Samples.
arXiv:2303.08559, 2023.

\bibitem{lewis2020rag}
Lewis P. et al.
Retrieval-Augmented Generation for Knowledge-Intensive NLP Tasks.
Advances in Neural Information Processing Systems, 2020.

\end{thebibliography}

\appendix
\chapter{Шаблоны промптов}

Для воспроизводимости ниже приведены шаблоны промптов. В запусках каждый шаблон получал один термин и требовал вернуть только список через \texttt{;}.

\textbf{zero.}
\begin{verbatim}
Перечисли гиперонимы для русского существительного.
Ответ: только гиперонимы через "; " (точка с запятой и пробел).
Никаких пояснений.

гипоним: {WORD}
гиперонимы:
\end{verbatim}

\textbf{cot} (Chain-of-thought).
\begin{verbatim}
Сначала подумай, какие гиперонимы наиболее вероятны,
но не показывай ход рассуждений.
Затем выведи только итоговый список (максимум 15) через "; ".

гипоним: {WORD}
гиперонимы:
\end{verbatim}

\textbf{ontology.}
\begin{verbatim}
Перечисли гиперонимы для русского существительного так,
как их обычно задают в тезаурусе (краткие словарные наименования).

Требования:
- только существительные или именные группы;
- без пояснений, примеров и комментариев;
- от более общего к менее общему;
- избегай слишком частных и случайных ассоциаций;
- максимум 15;
- ответ только через "; " (точка с запятой и пробел).

гипоним: {WORD}
гиперонимы:
\end{verbatim}

\textbf{rkt\_basic.}
\begin{verbatim}
Перечисли гиперонимы для термина из области ракетно-космической техники.
Используй принятую в этой области терминологию (летательный аппарат,
ракета-носитель, космический аппарат, баллистическая ракета и т. п.,
а не бытовые слова "техника", "объект").
Ответ: только гиперонимы через "; ". Никаких пояснений.

термин: {WORD}
гиперонимы:
\end{verbatim}

\textbf{rkt\_few\_shot.}
\begin{verbatim}
Ты эксперт по тезаурусу ракетно-космической техники.
Дай гиперонимы термина в стиле этого тезауруса.
Примеры:
СПУТНИК СВЯЗИ        -> искусственный спутник Земли; космический аппарат
БАЛЛИСТИЧЕСКАЯ РАКЕТА -> ракета; вооружение
ВЕРТОЛЁТ МИ-8        -> вертолёт; летательный аппарат

Ответ: только гиперонимы через "; ". Максимум 15.

термин: {WORD}
гиперонимы:
\end{verbatim}

\textbf{rkt\_role.}
\begin{verbatim}
Ты составляешь предметный тезаурус ракетно-космической техники.
Для термина выпиши его непосредственные гиперонимы, как в академическом
тезаурусе:
- ближайший родовой концепт ставь первым,
- избегай слишком общих слов ("объект", "техника", "сущность"),
- избегай метафор и описаний,
- от 3 до 10 элементов через "; ".

термин: {WORD}
гиперонимы:
\end{verbatim}

\end{document}